\subsection{局部连通空间}

定义 4. 如果拓扑空间 X 的每一点都有一个由连通邻域组成的基本邻域系，就称 X 为局部连通的。

* 我们将在第四章，§ 2，第 5 目看到，实直线是局部连通空间。*

在空间 X 的每点 x 处存在一个连通邻域，绝不能推出 X 是局部连通的。特别地，X 可以连通而非局部连通（习题 2 与 13）。反之，空间可以局部连通而不连通（例如含有多于一个点的离散空间）。

命题 11. 空间 X 局部连通的必要充分条件是：X 中开集的每个分支在 X 中都是开的。

该条件是充分的，因为这时 x 相对于 x 的一个开邻域的分支是 x 在 X 中的一个邻域。

反之，设 A 是局部连通空间 X 的开子集，B 是 A 的一个分支，且设 \(x \in B\)。设 V 是包含于 A 中的 x 的一个连通邻域；由分支的定义，V 含于 B；故 B 在 X 中是开的（§ 1，第 2 目，命题 1）。

因此，局部连通空间 X 的诸分支构成 X 的一个分成开集的分拆，从而 X 是它的诸分支的和（§ 2，第 4 目，例 3）。

推论. 设 U 是局部连通空间 X 的开子集，V 是 U 的一个分支。则 V 的边界（相对于 X）含于 U 的边界。

因为 V 在 U 中既开又闭，故 V 的边界点（相对于 X）不能属于 U，否则它也将是 V 相对于 U 的边界点，而这样的点不存在。

命题 12. 局部连通空间的每个商空间都是局部连通的。

设 X 是局部连通空间，R 是 X 上的等价关系，\(\varphi: X \to X/R\) 是典范映射。设 A 是

X/R 的一个开子集，C 是 A 的一个分支。则 \(\overline{\varphi}^{1}(C)\) 是 \(\overline{\varphi}^{1}(A)\) 的诸分支之并；因为若 \(x \in \overline{\varphi}^{1}(C)\)，且 K 是 \(x\) 在 \(\overline{\varphi}^{1}(A)\) 中的分支，则 \(\varphi(K)\) 是连通的（第 2 目，命题 4），含于 A，且含有 \(\varphi(x)\)；于是由 C 的定义有 \(\varphi(K) \subset C\)，因此 \(K \subset \overline{\varphi}^{1}(C)\)。既然 X 局部连通而 \(\overline{\varphi}^{1}(A)\) 在 X 中是开的，由命题 11 得 \(\overline{\varphi}^{1}(C)\) 在 X 中是开的；从而 C 在 X/R 中是开的，故再由命题 11，X/R 是局部连通的。

命题 13. a) 设 \((\mathbf{X}_t)_{t \in I}\) 是局部连通空间的一个族，使得 \(\mathbf{X}_t\) 除有限多个指标 \(t \in I\) 外都是连通的。则积空间 \(\mathbf{X} = \prod_{i \in I} \mathbf{X}_t\) 是局部连通的。

\begin{enumerate}
\def\labelenumi{\alph{enumi})}
\setcounter{enumi}{1}
\item
  反之，若非空拓扑空间的族 \((\mathbf{X}_{t})\) 之积是局部连通的，则每个 \(X_{t}\) 都局部连通，且 \(X_{t}\) 除有限多个指标外都连通。
\item
  设 \(\mathbf{J}\) 是 \(\mathbf{I}\) 的有限子集，使得 \(\mathbf{X}_i\) 不连通当且仅当 \(i \in \mathbf{J}\)。设
\end{enumerate}

\[
\mathrm{U} = \prod_ {i \in I} \mathrm{U} _ {i}
\]

是含有一点 \(x = (x_{t})\)（属于 \(\mathbf{X}\)）的初等集，并设 \(\mathbf{K}\) 是 \(\mathbf{I}\) 的有限子集，使得 \(\mathbf{U}_{t} \neq \mathbf{X}_{t}\) 当且仅当 \(t \in \mathbf{K}\)。令 \(\mathbf{V}_{t}\) 为 \(\mathbf{X}_{t}\)（当 \(t \notin J \cup K\)），而令 \(V_{t}\) 为 \(x_{t}\) 的一个含于 \(U_{t}\) 中的连通邻域（当 \(t \in J \cup K\)）；则

\[
\mathbf {V} = \prod_ {i \in I} \mathbf {V} _ {i}
\]

是连通的（由第 4 目命题 8），且是包含于 U 中的 x 的一个邻域。故 X 局部连通。

\begin{enumerate}
\def\labelenumi{\alph{enumi})}
\setcounter{enumi}{1}
\tightlist
\item
  设 \(a = (a_{i})\) 是 \(\mathbf{X}\) 的一点，\(\mathbf{V}\) 是 \(a\) 在 \(\mathbf{X}\) 中的一个连通邻域。由于除有限多个指标外有 \(\mathrm{pr}_{\mathfrak{i}}\mathbf{V} = \mathbf{X}_{\mathfrak{i}}\)（§ 4，第 1 目），由第 2 目命题 4，诸 \(\mathbf{X}_{\mathfrak{i}}\) 除有限多个指标外都是连通的。另一方面，对每个 \(x \in I\)、每个 \(a_{x} \in X_{x}\) 以及邻域 \(V_{x}\)（即 \(a_{x}\) 在 \(X_{x}\) 中的邻域），存在一点 \(x\)（属于 \(\mathbf{X}\)）使得 \(\mathrm{pr}_{\mathfrak{x}}x = a_{x}\)，且
\end{enumerate}

\[
\mathbf {V} = \mathbf {V} _ {x} \times \prod_ {i \neq x} \mathbf {X} _ {i}
\]

是 \(x\) 在 \(\mathbf{X}\) 中的一个邻域；于是 \(\mathbf{V}\) 含有一个连通邻域 \(W\)（即 \(x\) 的连通邻域），其投影 \(\mathrm{pr}_xW\) 是 \(a_x\) 的一个连通邻域，且含于 \(V_x\) 中（第 2 目，命题 4 及 § 4，第 2 目，命题 5）。故每个 \(X_x\) 都局部连通。

